\chapter{Chiral diagonals and curvature}
\label{nd:chap:diagonal}

Cohomological joining uses pullback and Gysin maps of a target diagonal.
An operation on chiral states must also account for the momentum fields
conjugate to its normal directions. The simplest operator
product prevents restriction by a quotient: killing a coordinate would
kill the vacuum. A differential can instead remove the coordinate and
its conjugates together.

We construct this differential first in an affine chart and then on a
normal bundle. Its cohomology retains the tangential chiral de Rham
fields. Passing from the normal bundle to the actual formal diagonal
introduces curvature. A relative vertex complex retains that formal
geometry and its fermionic charge, while a second complex, obtained by
finite conjugation, resolves the central Dolbeault algebra. The final
calculation records the parameter torsion that appears when the normal
contraction is made compatible with Laurent charge. These constructions
specify what a chiral version of the cochain maps in
Chapter~\ref{df:chap:fibers} would have to preserve.

\section{The normal quartet}
\label{nd:sec:quartet}

Work over $\mathbb C$. In one normal coordinate the chiral de Rham
system has even fields $v,\beta$ and odd fields $c,b$, with singular
operator products
\begin{equation}\label{nd:eq:ope}
 \beta(z)v(w)\sim\frac1{z-w},\qquad
 b(z)c(w)\sim\frac1{z-w}.
\end{equation}
The variables $z,w$ are insertion variables, distinct from target
coordinates such as $v$.

All other generating singular products vanish,
and vertex skew symmetry determines the reversed products. In
particular $v(z)\beta(w)\sim-(z-w)^{-1}$.

We use the following weights and gradings:
\[
\begin{array}{c|rrrr}
 &v&\beta&c&b\\\hline
 \text{conformal weight}&0&1&0&1\\
 \text{fermionic charge }J&0&0&1&-1\\
 \text{cochain degree}&0&0&-1&1.
\end{array}
\]
For these fields, cochain degree is $-J$ and its reduction modulo two
is their parity. These are not two independent contributions to a
tensor-product sign. Later, Dolbeault form degree will be added to
cochain degree.

The vertex products $a_{(n)}b$ give the singular part
\[
 a(z)b(w)\sim\sum_{n\geq0}
           \frac{(a_{(n)}b)(w)}{(z-w)^{n+1}}.
\]
Colons denote normal ordering, and $T$ denotes vertex translation.

The local systems \eqref{nd:eq:ope}, with the chiral coordinate-change
laws, define the sheaf $\Omega_Y^{\mathrm{ch}}$ of vertex
superalgebras on a smooth complex manifold $Y$.\footnote{
F.~Malikov, V.~Schechtman and A.~Vaintrob,
\emph{Chiral de Rham complex},
\href{https://arxiv.org/abs/math/9803041}{Sections~2--3}.
The coordinate transformations and their composition are in
Sections~3.6--3.9.}
Here it has zero differential. In particular, its usual chiral
de Rham differential is not included.

Suppose that a unital vertex homomorphism $\phi$ kills the normal
coordinate. The relation $\beta_{(0)}v=1$ would give
\[
 1=\phi(\beta_{(0)}v)=\phi(\beta)_{(0)}\phi(v)=0.
\]
Thus a chiral restriction cannot be this quotient.

Define instead
\begin{equation}\label{nd:eq:delta}
 \delta=(:vb:)_{(0)},\qquad
 \delta c=v,\quad \delta\beta=-b,\quad
 \delta v=\delta b=0.
\end{equation}
Every zero mode is a derivation of every vertex product and commutes
with $T$. The current $:vb:$ has regular self-product, hence
$\delta^2=0$. The sign on the momentum field is forced by
\[
 \delta(\beta_{(0)}c)=(-b)_{(0)}c+\beta_{(0)}v=-1+1=0.
\]
The operator $\delta$ raises cochain degree by one, preserves conformal
weight, and lowers $J$ by one.

We use two carriers. The polynomial carrier consists of finite normally
ordered states with polynomial normal coordinates.

For the formal carrier,
complete the coefficient functions in the undifferentiated normal
coordinates at each fixed conformal weight. Then take the direct sum
over conformal weights. A state consequently has bounded conformal
weight. Positive-weight creation variables remain polynomial, with
bounded multiplicities in each weight. Singular products differentiate
coefficient functions only finitely many times for a fixed pair of
states, and are continuous in this topology.

\begin{theorem}\label{nd:thm:quartet}
Tensor a tangential chiral de Rham algebra with finitely many normal
quartets \eqref{nd:eq:ope}. Give the tensor product the sum of the
normal differentials \eqref{nd:eq:delta}, and give the tangential
factor zero differential and cochain grading $-J$.
The tangential inclusion is a
quasi-isomorphism of differential graded vertex superalgebras on both
the polynomial and the fixed-weight formal carriers.
\end{theorem}

\begin{proof}
Use the Poincar\'e--Birkhoff--Witt (PBW) state presentation, whose
underlying vector space has polynomial and exterior creation variables.
Write $v_j,\beta_j,c_j,b_j$ for the creation variables corresponding
to the normalized derivatives $T^j/j!$ of the four fields, $j\geq0$.

On one normal factor,
\[
 \delta=\sum_{j\geq0}
       (v_j\partial_{c_j}-b_j\partial_{\beta_j}),\qquad
 h_0=\sum_{j\geq0}
       (c_j\partial_{v_j}-\beta_j\partial_{b_j}).
\]
Fermionic derivatives are left derivatives. The anticommutator
$\delta h_0+h_0\delta$ is a derivation of the PBW polynomial and
exterior algebra and is the identity on every generator. It is
therefore the total normal-number operator $N$.
For several normal factors, sum the two operators over the factors.

On positive normal number put $h=h_0/N$, and put $h=0$ on the
normal vacuum. Let $i$ include that vacuum and let $\pi$ extract
its coefficient. Then
\begin{equation}\label{nd:eq:contraction}
 \delta h+h\delta=1-i\pi.
\end{equation}
Each polynomial involves finitely many modes.

On formal coefficient
functions the same identity holds coefficientwise. Differentiation
lowers adic order by at most one. Division by a positive integer
preserves the topology. Thus the normal cohomology is its scalar
vacuum.

Tensoring this contraction with the tangential factor proves
the cohomology assertion. The inclusion preserves the vacuum,
translation and all vertex products, so it is the asserted vertex
quasi-isomorphism.
\end{proof}

The projection $\pi$ is a map of complexes. The equation
$\beta_{(0)}v=1$ prevents it from being a vertex homomorphism.
The homotopy is likewise a linear operator, not a vertex derivation.
These distinctions allow the contraction to exist without contradicting
the obstruction to a quotient.

The weight bound is essential for locality. An unrestricted oscillator
completion would admit both $\sum_{m\geq0}:\beta^m:$ and
$\sum_{m\geq0}:v^m:$. Their complete contractions include
$m!(z-w)^{-m}$ for arbitrarily large $m$, so no finite locality order
would remain. The first series has unbounded conformal weight and is
excluded.

Finally, $\delta$ does not preserve an individual Laurent coefficient
of the elliptic state, because it lowers fermionic charge. It removes
only normal fields. Tangential states of positive conformal weight
survive. A differential contracting all positive conformal weights
has a different cohomology.

\section{Affine diagonals and regrouping}
\label{nd:sec:affine}

Consider $r\geq1$ copies of one affine target coordinate, with field
columns $\gamma,c,\beta,b$. Let $\mathbf1$ be the column of $r$ ones
and set
\begin{equation}\label{nd:eq:projector}
 P_r=I_r-\frac1r\mathbf1\mathbf1^{\mathsf t},\qquad
 \delta_r=
 \left(:\sum_{i,j}(P_r)_{ij}\gamma_i b_j:\right)_{(0)}.
\end{equation}
The fields tangent to the full diagonal are
\[
 \bar\gamma=\frac1r\sum_i\gamma_i,\qquad
 \bar c=\frac1r\sum_i c_i,\qquad
 \bar\beta=\sum_i\beta_i,\qquad
 \bar b=\sum_i b_i.
\]

Their two nonzero generating products have coefficient one, as in
\eqref{nd:eq:ope}. The complementary fields form $r-1$ normal
quartets. On the original columns,
$\delta_rc=P_r\gamma$ and $\delta_r\beta=-P_rb$.
Theorem~\ref{nd:thm:quartet} contracts their total normal number.

\begin{proposition}\label{nd:prop:regrouping}
Successive affine diagonal operations associated to any partition
and refinement of the $r$ factors give the differential $\delta_r$.
Their changes of normal basis compose and preserve the tangential
fields and all vertex products.
\end{proposition}

\begin{proof}
For a set partition $\mathcal P$, let $Q_{ij}=1/|B|$ when $i,j$
belong to the same block $B\in\mathcal P$, and let $Q_{ij}=0$
otherwise. It is the orthogonal projector onto block-constant vectors.
There is an orthogonal decomposition
\[
 P_r=(I_r-Q)+\left(Q-\frac1r\mathbf1\mathbf1^{\mathsf t}\right).
\]

The first term records the normal directions within each block.
The second records the differences of block means, with the block-size
weights in their induced pairing. Their currents sum to the current
in \eqref{nd:eq:projector}. Repeated refinement decomposes the same
projector.

A normal-basis change sends $v,c$ by an invertible matrix and
$\beta,b$ by its inverse transpose. It preserves the two pairings
in \eqref{nd:eq:ope} and the current $\sum_a:v^ab_a:$.
Composition of these transformations is matrix multiplication, proving
the compatibility of all successive regroupings.
\end{proof}

This calculation concerns the diagonal in a product of affine target
charts. To apply it to a curved target, one must first understand its
behavior under changes of normal frame.

\section{The absolute normal-bundle complex}
\label{nd:sec:absolute}

Let $E\to Y$ be a holomorphic vector bundle, and let
$M=\operatorname{Tot}(E)$. Locally write $x^i$ for base coordinates
and $v^a$ for fiber coordinates. There is a chiral de Rham quartet
for every base and fiber coordinate. Let $\mathcal V_E$ be the sheaf
over $Y$ of states on $M$ polynomial in the fiber directions, or its
fixed-weight completion along the zero section.

Give every $c$, including each $c_x^i$, degree $-1$, and every $b$
degree $1$. This total fermionic grading respects a varying frame:
under $v'=g(x)v$ one has
\begin{equation}\label{nd:eq:odd-frame}
 c_v'=g(x)c_v+(dg)_x(v)c_x.
\end{equation}
A degree assigned only to normal fermions would not respect
\eqref{nd:eq:odd-frame}.

The fiber Euler vector field $\mathsf E=v^a\partial_{v^a}$ is
intrinsic. Chiral contraction with a vector field is compatible with
the chiral de Rham coordinate laws. Consequently
\begin{equation}\label{nd:eq:absolute-delta}
 \delta=(\iota_{\mathsf E})_{(0)},\qquad
 \iota_{\mathsf E}=\sum_a:v^ab_a:
\end{equation}
is a global differential on $\mathcal V_E$.
In a local frame it has the normal action
\eqref{nd:eq:delta} and kills tangential fields.

The following homotopy controls the nonconstant part of a frame change.
For a matrix $A(x)$, write $W=A^b_a(x)v^a\partial_{v^b}$ and put
\begin{equation}\label{nd:eq:frame-homotopy}
 H_A=:A^b_a c_v^a\beta_b:
       -:\partial_iA^b_a\,c_x^ic_v^ab_b:.
\end{equation}
The chiral Lie derivative along $W$, denoted by $\mathcal L_W$,
is the infinitesimal action of its coordinate change.

\begin{lemma}\label{nd:lem:frame}
The zero mode of \eqref{nd:eq:frame-homotopy} satisfies
$[\delta,(H_A)_{(0)}]=\mathcal L_W$.
\end{lemma}

\begin{proof}
Applying the odd derivation $\delta$ gives the current
\[
 :A^b_av^a\beta_b:
 +:A^b_ac_v^ab_b:
 +:\partial_iA^b_a\,v^ac_x^ib_b:.
\]
The last sign follows from
$\delta(c_x^ic_v^ab_b)=-c_x^iv^ab_b$.
The three terms are the vector-field term and the two fermionic terms
in the chiral Lie derivative of $W$. Taking zero modes proves the
identity.
\end{proof}

\begin{theorem}\label{nd:thm:absolute}
The cohomology vertex sheaf of $(\mathcal V_E,\delta)$ is canonically
$\Omega_Y^{\mathrm{ch}}$, with zero differential and cochain grading
equal to minus its fermionic charge. This holds for polynomial normal
states and for the fixed-weight formal completion.
\end{theorem}

\begin{proof}
In each frame, Theorem~\ref{nd:thm:quartet} identifies the cohomology
with the tangential fields. It remains to compare these identifications.
After shrinking around a base point, a varying frame matrix is a
constant matrix times $\exp A(x)$, using the holomorphic matrix
logarithm near the identity. A constant frame matrix fixes the
tangential fields and the normal vacuum.

Let $F_s$ be the chiral frame action of $\exp(sA)$. It commutes with
$\delta$ because $[\mathsf E,W]=0$, and
$\frac{d}{ds}F_s=\mathcal L_WF_s$.
Lemma~\ref{nd:lem:frame} gives
\[
 F_1-1=\delta K+K\delta,
 \qquad K=\int_0^1(H_A)_{(0)}F_s\,ds.
\]
Thus every vertical frame change induces the identity on the local
tangential cohomology. A change of base coordinates induces the base
chiral de Rham transformation. These identifications satisfy the
cocycle identity and preserve vertex products, giving the canonical
cohomology sheaf.
\end{proof}

The theorem constructs a canonical cohomology sheaf without requiring
a global strict inclusion of tangential vertex fields.
For the full diagonal in $Y^r$, its linear normal bundle is
$T_Y\otimes(\mathbb C^r/\mathbb C_{\mathrm{diag}})$.
The theorem applies to this bundle, and the decompositions used in
Proposition~\ref{nd:prop:regrouping} give its compatible linear
regroupings. The next issue is comparison with the formal diagonal
itself.

\section{Curvature of the formal diagonal}
\label{nd:sec:curvature}

Write two nearby target points as $(x,x+v)$. Under $x'=f(x)$,
their difference becomes
\begin{equation}\label{nd:eq:hessian}
 v'=f(x+v)-f(x)
    =df_x(v)+\tfrac12d^2f_x(v,v)+O(v^3).
\end{equation}
The quadratic term is the first discrepancy from the linear normal
transition.

\begin{proposition}\label{nd:prop:connection}
Let $Y$ be a K3 or Enriques surface. Its formal diagonal has no
holomorphic linearization over the first projection inducing the
identity on the normal bundle.
\end{proposition}

\begin{proof}
The quadratic term of such a linearization, written as an exponential,
would have the form
\[
 \operatorname{Exp}_x(v)=x+v-\tfrac12\Gamma_x(v,v)+O(v^3)
\]
with $\Gamma_x$ symmetric. Comparing
$f(\operatorname{Exp}_x(v))$ and
$\operatorname{Exp}_{f(x)}(df_x(v))$ gives
\[
 \Gamma_{f(x)}(df_x(v),df_x(v))
       =df_x\Gamma_x(v,v)-d^2f_x(v,v).
\]
This is the transformation law of a torsion-free holomorphic
connection on $T_Y$.

Extend that connection by the bundle's $\bar\partial$ operator.
Its curvature has type $(2,0)$, so the Chern--Weil form representing
$c_2$ has type $(4,0)$ and vanishes on a complex surface.
This would give $c_2(T_Y)=0$, contrary to
$\int_Yc_2(T_Y)=24$ for a K3 surface and $12$ for an Enriques surface.
\end{proof}

Equivalently, the quadratic discrepancies in
\eqref{nd:eq:hessian} represent the symmetric Atiyah class of $T_Y$,
which obstructs a holomorphic connection. This obstruction leaves the
ordinary diagonal kernel intact. The sheaf $\mathcal O_\Delta$ is
perfect because the local coordinate differences are a regular
sequence. For a unital affine coordinate algebra $B$, the bar complex has terms
$B^{\otimes(n+2)}$, outer bimodule action and alternating adjacent
multiplication. It resolves $B$ as a $B$-bimodule functorially.
A chiral comparison must retain the additional coordinate-change data
of its normal fields.

\section{A relative lift retaining the formal neighborhood}
\label{nd:sec:formal}

Assume that $Y$ is Kähler. Choose a Kähler connection and its smooth
formal exponential. The transported Dolbeault algebra of the formal
diagonal over the first projection is
\begin{equation}\label{nd:eq:formal-base}
 \left(\mathcal A_Y^{0,*}
       (\widehat{\operatorname{Sym}}T_Y^*),D_B\right),
 \qquad D_B=\bar\partial+R,\qquad R=\sum_{n\geq2}R_n.
\end{equation}
Here $\mathcal A_Y^{0,*}$ is the Dolbeault algebra of smooth
$(0,*)$-forms. Each $R_n$ is a $(0,1)$-form with values in
homogeneous vertical polynomial vector fields of degree $n$.

We define $R_n$ as the homogeneous Taylor coefficient of this
transported differential. This fixes the factorial convention.

The formal-exponential theorem gives
\begin{equation}\label{nd:eq:MC}
 \bar\partial R+\tfrac12[R,R]=0,
\end{equation}
and $R_2$ represents the Atiyah class.\footnote{
S.~Yu, \emph{The Dolbeault dga of the formal neighborhood of the
diagonal}, \href{https://arxiv.org/abs/1211.1567}{Theorem~3.2 and
Section~3}. The homogeneous polynomial convention in
\eqref{nd:eq:formal-base} is obtained from its symmetric-tensor
Taylor convention.}

Define the relative formal $bc$--$\beta\gamma$ system
$\mathcal V_{\mathrm{rel}}$ of $T_Y$ as follows. Its normal fields
have products \eqref{nd:eq:ope}, its coefficients
$\mathcal A_Y^{0,*}$ are central vertex elements with $T=0$, and
it has no base momentum fields. Linear frames act on $v,c$ and
contragrediently on $\beta,b$. Complete in the normal coordinates
at each conformal weight, as in Section~\ref{nd:sec:quartet}.

Total cochain degree is Dolbeault degree minus $J$. Forms are placed
to the left of the fiber fields. The tensor product uses total
degree for its signs.

For a vertical vector field $V=V^i(v)\partial_{v^i}$, its chiral
Lie action is
\begin{equation}\label{nd:eq:lie}
 \mathcal L_V=
 \left(:V^i(v)\beta_i:
       +:\partial_jV^i(v)c^jb_i:\right)_{(0)}.
\end{equation}
It acts by $V$ on functions and by
$c^i\mapsto\partial_jV^i c^j$ on the coordinate fermions.
This action is the derivative of the formal-coordinate action in
the chiral de Rham construction. Its bosonic and fermionic anomalies
cancel, giving the exact identity
\begin{equation}\label{nd:eq:lie-action}
 [\mathcal L_V,\mathcal L_W]=\mathcal L_{[V,W]}.
\end{equation}
We extend it to form-valued vector fields with the tensor-product signs.

\begin{theorem}\label{nd:thm:formal}
The operator
\begin{equation}\label{nd:eq:Dform}
 D_{\mathrm{form}}=\bar\partial+\mathcal L_R
\end{equation}
is a global square-zero derivation of $\mathcal V_{\mathrm{rel}}$.
It preserves conformal weight and fermionic charge, and restricts to
$D_B$ on fiber functions. Changes of formal exponential induce chiral
coordinate isomorphisms intertwining these differentials, with the
coordinate composition law.
\end{theorem}

\begin{proof}
Equations \eqref{nd:eq:MC} and \eqref{nd:eq:lie-action} give
\[
 D_{\mathrm{form}}^2
   =\mathcal L_{\bar\partial R+[R,R]/2}=0.
\]
The current in \eqref{nd:eq:lie} has weight one and charge zero.
Its zero mode is a derivation, preserves these gradings, and acts by
$V$ on functions.

Under a formal coordinate change, the chiral de Rham
transformation transports the Lie action. For a base-dependent change,
its Dolbeault derivative supplies the extra term that also appears
when conjugating $\bar\partial+R$ on functions. Hence the chiral
transformation intertwines \eqref{nd:eq:Dform}. Composition of the
coordinate transformations gives the cocycle identity.

At fixed conformal weight, a state has boundedly many momentum fields.
Only boundedly many contractions can therefore lower the normal order
of a coefficient. The sum of homogeneous terms is continuous in the
specified formal completion. Thus these local formulas define the
global formal operator.
\end{proof}

The relative Euler contraction $\delta$ of
\eqref{nd:eq:absolute-delta} is still defined on this carrier. Its
commutator with the new differential measures the curvature error.
For a form-valued vertical field $V$, write
$\iota_V=:V^ib_i:$.

\begin{proposition}\label{nd:prop:curvature}
With the total graded commutator,
\begin{equation}\label{nd:eq:curvature}
 [D_{\mathrm{form}},\delta]
   =-\sum_{n\geq2}(n-1)(\iota_{R_n})_{(0)}.
\end{equation}
In particular, $D_{\mathrm{form}}+\delta$ need not square to zero.
\end{proposition}

\begin{proof}
The chiral Cartan identity gives
$[\mathcal L_V,\iota_{\mathsf E,(0)}]
=\iota_{[V,\mathsf E],(0)}$.
For a homogeneous degree-$n$ vector field,
$[V,\mathsf E]=-(n-1)V$.
The holomorphic Euler contraction commutes with $\bar\partial$
with the total sign. Extending the Cartan identity to forms yields
\eqref{nd:eq:curvature}.
\end{proof}

\section{A conjugated Koszul resolution}
\label{nd:sec:koszul}

The Euler contraction can also be used to construct a global
resolution of the central base. For a degree-$n$ vertical field put
\begin{equation}\label{nd:eq:H}
 H_{V_n}=\frac1n:\partial_jV_n^i(v)c^j\beta_i:.
\end{equation}
The odd derivation rule and the homogeneous Euler identity give
\begin{equation}\label{nd:eq:K}
 \delta H_{V_n}=K_{V_n}
  =:V_n^i\beta_i:
       +\frac1n:\partial_jV_n^ic^jb_i:.
\end{equation}
The fermionic coefficient here is $1/n$, whereas the Lie action
\eqref{nd:eq:lie} has coefficient one.

Apply \eqref{nd:eq:H} to the tensors $R_n$ and let
$\mathscr S=\sum_n H_{R_n}$, with forms placed first. Its form degree
is one and its fiber cochain degree is $-1$.
Thus $S_0=\mathscr S_{(0)}$ is an even derivation of total degree zero.
Define
\begin{equation}\label{nd:eq:DK}
 d_0=\bar\partial+\delta,\qquad
 D_K=e^{S_0}d_0e^{-S_0}.
\end{equation}
On a complex $d$-fold, $S_0^{d+1}=0$ by form degree, so both
exponentials are finite.

\begin{theorem}\label{nd:thm:koszul}
The operator $D_K$ is a global square-zero vertex derivation.
The central inclusion
\[
 i:(\mathcal A_Y^{0,*},\bar\partial)
        \longrightarrow(\mathcal V_{\mathrm{rel}},D_K)
\]
is a quasi-isomorphism. For choices $S_0,S'_0$ on this carrier,
the automorphism $e^{S'_0}e^{-S_0}$ intertwines their resolutions.
These intertwiners compose exactly.
\end{theorem}

\begin{proof}
The currents \eqref{nd:eq:H} are tensorial under relative linear
frame changes. Their zero modes are global derivations, and finite
exponentiation gives vertex automorphisms fixing the central base.
Conjugation therefore proves the square-zero and derivation assertions.

The normal homotopy \eqref{nd:eq:contraction} is equivariant under
linear frame changes with central coefficients. It defines a global
relative homotopy $h$ and vacuum projection $\pi$. Extend $h$ to
forms with the tensor-product sign. Then
$\bar\partial h+h\bar\partial=0$ and
$d_0h+hd_0=1-i\pi$.

Transporting the contraction gives
\[
 h_K=e^{S_0}he^{-S_0},\qquad
 \pi_K=\pi e^{-S_0},\qquad
 D_Kh_K+h_KD_K=1-i\pi_K.
\]
This proves the quasi-isomorphism of complexes, and the central
inclusion is a vertex morphism. The projection and the homotopy are
only required to be linear maps of the indicated degrees.
Finally, cancellation of adjacent exponentials proves both the
intertwining identity and composition of the changes of choice.
\end{proof}

The square-zero argument uses homogeneous form-valued tensors to
define $\mathscr S$, but does not require their Maurer--Cartan equation.
That equation is needed for the geometric interpretation of
$D_{\mathrm{form}}$, which restricts to the formal-neighborhood
Dolbeault differential on functions. Theorem~\ref{nd:thm:koszul}
compares $(\mathcal V_{\mathrm{rel}},D_K)$ with the central Dolbeault
algebra by a quasi-isomorphism. It does not identify the two differentials.

\section{The second correction on a surface}
\label{nd:sec:surface}

On a surface the conjugation has only two nontrivial form degrees.
Write $K_1=\sum_nK_{R_n}$ as a form-valued current. The total sign is
$d_0\mathscr S=\bar\partial\mathscr S-K_1$.
Expanding \eqref{nd:eq:DK} gives
\begin{equation}\label{nd:eq:surface}
 D_K=d_0+(K_1)_{(0)}-(\bar\partial\mathscr S)_{(0)}
                    +\tfrac12[S_0,(K_1)_{(0)}].
\end{equation}
All further terms have form degree at least three.

The first-form action on functions is $R$. The second-form term
can contain fermionic fields of negative cochain degree.

\begin{example}\label{nd:ex:second}
Take fiber coordinates $v_1,v_2$ and exterior form variables
$\eta_1,\eta_2$, placed to the left of fiber fermions. Set
\[
 V=v_1^2\partial_{v_2},\qquad W=v_2^2\partial_{v_1},\qquad
 H_V=:v_1c_1\beta_2:,\qquad H_W=:v_2c_2\beta_1:,
\]
and suppress $\bar\partial$. For
$S_0=\eta_1(H_V)_{(0)}+\eta_2(H_W)_{(0)}$, put
$A_1=[S_0,\delta]$ and $A_2=\tfrac12[S_0,A_1]$.
Then
\[
 A_1(v_1)=\eta_2v_2^2,\qquad
 A_1(v_2)=\eta_1v_1^2,
\]
while
\begin{align*}
 A_2(v_1)&=-\eta_1\eta_2
       \left(\tfrac32v_1v_2c_1+\tfrac12v_1^2c_2\right),\\
 A_2(v_2)&=\eta_1\eta_2
       \left(\tfrac12v_2^2c_1+\tfrac32v_1v_2c_2\right).
\end{align*}

Indeed $S_0(v_1)=\eta_2v_2c_2$, and the two compositions are
\[
 S_0A_1(v_1)=-2\eta_1\eta_2v_1v_2c_1,\qquad
 A_1S_0(v_1)=\eta_1\eta_2(v_1v_2c_1+v_1^2c_2).
\]
Their half-difference gives the first formula. Interchanging the two
indices gives the second. Each action on these coordinate polynomials
has only a single contraction, so these are identities of zero modes.
\end{example}

Thus the algebra of fiber functions need not be a subcomplex of
$D_K$. The example uses specified polynomial tensors. They are not
assumed to satisfy \eqref{nd:eq:MC} or to arise from a K3 metric.
They show that the second correction can be nonzero for arbitrary
tensors.

\section{Charge restoration and parameter torsion}
\label{nd:sec:charge}

Introduce an even parameter $u$ of fermionic charge $+1$ and
cochain degree zero. The combined charge includes its exponent.
Cochain degree still uses the original field grading and Dolbeault
degree. In particular, charge is no longer parity after adjoining $u$.
The differential $u\delta$ preserves combined charge.

Let $K$ be a finite normal PBW subcomplex stable under $\delta$ and
the contraction $h$, with cohomology consisting of its scalar vacuum
in degree zero. Write
$B^k=\operatorname{im}(\delta:K^{k-1}\to K^k)$.

\begin{proposition}\label{nd:prop:torsion}
There are vector-space splittings giving
\begin{equation}\label{nd:eq:torsion}
 H^k(K[u],u\delta)\cong
 B^k\otimes\mathbb C[u]/(u)\ \oplus\
 \begin{cases}
 \mathbb C[u]|0\rangle,&k=0,\\
 0,&k\ne0.
 \end{cases}
\end{equation}
The derived fiber at $u=0$ is $(K,0)$.
The same statements hold coefficientwise in the fixed-weight formal
completion, with coefficientwise completed base change.
\end{proposition}

\begin{proof}
Multiplication by $u$ is injective on $K[u]$, so the cycles of
$u\delta$ are $(\ker\delta)[u]$ and its boundaries are $uB[u]$.
Split $\ker\delta$ into $B$ and the scalar vacuum to obtain
\eqref{nd:eq:torsion}. The finite complex is termwise free over
$\mathbb C[u]$, so derived base change sets its differential to zero.
In the fixed-weight completion the normal-number contraction gives
the same continuous splitting coefficientwise. Completed base change
has the asserted zero differential on each coefficient.
\end{proof}

For two normal directions, keep no derivative jets and truncate total
normal number at three. There are four bosonic and four fermionic
variables. Their normal-number series and its first four coefficients are
\[
 \frac{(1+t)^4}{(1-t)^4},\qquad 1+8+32+88=129.
\]
The contraction leaves one vacuum and $64$ acyclic pairs.
After scaling the differential there are therefore $64$ torsion
cohomology generators, besides the free vacuum. The terms of the
parameter complex are free even though its cohomology is not flat.

For the curved resolution, charge restoration begins over
$\mathbb C[u,u^{-1}]$. The finite conjugation is
\begin{equation}\label{nd:eq:homogenized}
 e^{S_0/u}(\bar\partial+u\delta)e^{-S_0/u}
  =\bar\partial+u\delta+(K_1)_{(0)}+u^{-1}K_2+\cdots.
\end{equation}
On a surface the last term is
\[
 K_2=-(\bar\partial\mathscr S)_{(0)}
                   +\tfrac12[S_0,(K_1)_{(0)}].
\]

A correction of form degree $r$ has charge $r-1$, so its coefficient
is $u^{1-r}$. A regular extension on the unchanged carrier requires
these negative-power coefficients to vanish, or to acquire the required
divisibility in a parameter-dependent family.
Example~\ref{nd:ex:second} exhibits a nonzero $K_2$ for polynomial
tensors. Its value for a specified K3 curvature is a separate
geometric calculation.

\section{The diagonal transfer still to be constructed}
\label{nd:sec:comparison}

For a compact complex surface $Y$ and $N\geq0$, the geometric state
object at length $N$ is
\begin{equation}\label{nd:eq:states}
 \mathcal H_N(Y)=
 \bigoplus_{[g]\subset\mathfrak S_N}
 R\Gamma\bigl(Y^N,\Omega_{Y^N}^{\mathrm{ch},g}\bigr)^{C(g)}.
\end{equation}
Here $C(g)$ is the centralizer of $g$. The sheaf
$\Omega_{Y^N}^{\mathrm{ch},g}$ consists of twisted chiral de Rham
modules on the fixed locus $(Y^N)^g$. We view it as a sheaf on $Y^N$.

Its coordinate and centralizer actions are those of the
twisted-module construction.\footnote{
E.~Frenkel and M.~Szczesny, \emph{Chiral de Rham complex and orbifolds},
\href{https://arxiv.org/abs/math/0307181}{Sections~4.8--4.10,
Theorems~4.2--4.3}. The index is treated in Section~6.}

The derived global sections in \eqref{nd:eq:states} use only the
sheaf-resolution differential. Finite-group invariants are exact
over $\mathbb C$. We use individual length and conformal-weight
components. Their generating series are taken coefficientwise.

The absolute Euler complex has cohomology $\Omega_Y^{\mathrm{ch}}$.
The relative formal-diagonal complex retains the transported
Dolbeault algebra on functions and the normal fields.
The conjugated Koszul complex resolves the central base.
Each construction therefore has a specified use in a comparison.
The supported diagonal transfer on
\eqref{nd:eq:states} remains to be constructed.

Let $S$ be a K3 surface and $E$ an elliptic curve. For the threefold
$S\times E$, the zero object in the relative translation quotient
lies over $BE$, the classifying stack of $E$.
The comparison must retain that relative unit and compatible
orientation square roots along extension flags.

It must also identify the reduced state with a specified geometric
module coefficient. Let $H,R$ be unital algebras and
$I_0:H\to R$ a unital homomorphism. A reduced functional
$D:H\to R$ satisfying
\[
 D(ab)=I_0(a)D(b)+D(a)I_0(b),\qquad D(1)=0
\]
defines the unital action
\[
 a\longmapsto\begin{pmatrix}I_0(a)&D(a)\\0&I_0(a)\end{pmatrix}
\]
on the right $R$-module $R\oplus R$. Multiplicativity is exactly the
displayed derivation identity. Thus $D$ is its off-diagonal coefficient. Its unit data differ from an
empty-vacuum matrix coefficient of value one.

A chiral McKay comparison with the Hilbert-scheme state complex
$R\Gamma(Y^{[N]},\Omega_{Y^{[N]}}^{\mathrm{ch}})$ must likewise identify
the complexes and their operations beyond their equal elliptic indices.

\begin{question}\label{nd:qn:transfer}
Construct a diagonal transfer on the actual chiral state sheaves,
through defects or bimodules preserving the prescribed Laurent charge.
Construct the compact reduced critical Hall coefficient functor,
with its oriented extension pullbacks and supported pushforwards,
and compare its diagonal operation with this transfer.
Can the comparison preserve the curvature corrections, derived
specialization, supported products, vacuum, and the common elliptic
torsor of the reduction?
\end{question}

The next coherence can be stated once its supports and coefficients
are fixed. Let $(B_A,b_A)$ and $(B_B,b_B)$ be specified supported bar
coalgebras with degree-one coderivations squaring to zero.
Suppose that the degree-zero components $F_1,\ldots,F_{n-1}$ satisfy
the comparison equations through arity $n-1$.
Let $\partial$ be the differential on the arity-$n$ Hom complex
induced by the unary coderivation components. The residual is
\[
 \mathfrak r_n=(b_BF_{<n}-F_{<n}b_A)_n,\qquad
 \partial F_n=-\mathfrak r_n.
\]
The preceding equations imply $\partial\mathfrak r_n=0$.

Suppose that this Hom complex contracts onto chosen cohomology
representatives, with projection $P$ and homotopy $h$. Then
$\partial h+h\partial=1-P$, so $F_n=-h\mathfrak r_n$ solves the
equation when $P\mathfrak r_n=0$.
The normal contraction in \eqref{nd:eq:contraction} computes its own
vertex complex. Applying this criterion to compact Hall operations
requires their supported Hom complex and the value of its residual.
